% Every simulated number is a macro from numbers.tex; claims.md gives its Kaggle notebook and file.
\subsection{Baseline Against Observations}
Over \ValSeeds{} seeds the baseline runs the Ring Road at \ValRing{} (survey 24.0) and inner main
roads at \ValInner{} (survey 15.3). It carries \ValShare{} of the counted volume summed over the
nine intersections, with a correlation of \ValCorr{} across them; on average \ValWithin{} of 9 are
within a factor of 1.5 and \ValGeh{} pass GEH below 5 (Figure~\ref{fig:val}). Measured the same
way, over the second hour of a two hour peak, the first version also carried about 64 percent of
the counts, but with a negative correlation across junctions ($-0.23$, one seed), and it
gridlocked, finishing 55 percent of trips; its published 46 percent used the first hour of an empty
network. The gain is in where traffic goes and in a stable network, not in total volume. The
remaining gap is concentrated at the four busiest junctions, Thapathali, Old and New Baneshwor and
Kalanki, which receive \ValBusyLo{} to \ValBusyHi{} percent of their counts. Halving or doubling
the assumed give way priority (0.1 and 0.4, five seeds each) moves inner speed by at most
\ValYieldKmh~km/h and the volume share by at most \ValYieldShare{} points.

\begin{figure}[ht]
\centering
\includegraphics[width=0.8\linewidth]{fig_jica_validation.png}
\caption{Model to observed peak entering volume at nine intersections (JICA 2011--12 counts), mean
over the baseline seed sweep (notebook 2). Error bars: 95\% CI.}
\label{fig:val}
\end{figure}

In the four hour horizon, \BaseComp{} of baseline trips finish, and mean time in system is
\BaseTis{} minutes. Volumes are 2011 levels; the network is a recent OpenStreetMap snapshot.

\subsection{Signal Control and Pricing}
Figure~\ref{fig:levers} reports all eight tests. Adaptive signal control changes time in system by
\HOneTis{} (\HOneNeg/15 seeds) and completed trip delay by \HOneDelay. Pricing changes time in
system by \HTwoTis{} (\HTwoNeg/15) and delay by \HTwoDelay. Both together change time in system by
\HThreeTis{} and delay by \HThreeDelay. All three survive the Bonferroni correction.

\begin{figure}[ht]
\centering
\includegraphics[width=0.8\linewidth]{fig_jica_levers.png}
\caption{Paired change in mean time in system, 15 seeds per configuration (notebooks 3a and 3b).
H4 is against morning demand without tidal lanes; the last row is against the volume matched
uniform cut. Grey: 95\% CI includes zero. Error bars: 95\% CI.}
\label{fig:levers}
\end{figure}

The combined lever is additive. Its per seed interaction term, the combined change minus the two
single changes, is \Inter{} of baseline time in system (\InterNeg{} of 15 seeds negative), centered
on zero. This agrees with the correction made to the first version, where an assumed super
additive effect did not survive a paired test.

\subsection{Capacity and Demand Management}
Staggered departures give the largest single effect, \HSixTis{} in time in system and \HSixDelay{}
in delay. Spreading the peak changes the share of trips finished at the four hour horizon by
\HSixComp, simply because they start later; time in system counts them, so the estimate is not
flattered by it. Both arms generate exactly \BaseVeh{} vehicles. A 50 percent longer peak is a
large coordination request in practice.

Decentralizing attraction changes time in system by \HSevenTis{} against the baseline, but cutting
every trip uniformly to the same total does as well: decentralized against uniform is
\HSevenUTis, not significant. In this model the benefit comes from fewer trips, not from where
they end. The first version reported a 10.5 percent advantage over the same control; it does not
reproduce here.

Tidal lanes do nothing measurable (\HFourTis). The parking enforcement proxy gives \HFiveTis{} in
time in system, which does not survive the Bonferroni correction, and no delay change
(\HFiveDelay). Both are read as null results.

\subsection{Comparison With the First Version}\label{sec:v1}
On the first version's synthetic demand and main road network, the same levers cut completed trip
delay by \OldAdaptive\% (adaptive), \OldPricing\% (pricing) and \OldComb\% (combined); these come
from local runs of that model, not the public run. Here the corresponding figures are
\HOneDelay, \HTwoDelay{} and \HThreeDelay. Adaptive signals keep, and here exceed, their effect;
pricing loses \OldPricingLoss{} of its effect and the combined lever \OldCombLoss. A likely reason is
route choice: tolls work by moving traffic to less used paths. With main roads only, parallel
arterials had spare capacity to absorb it; here the parallel options are side streets at a third
of a main road's capacity that must give way. This was not isolated by a separate experiment. The
first version loaded a one hour peak and this one a two hour peak, so only relative changes, not
absolute delays, are comparable between them.

\subsection{Demand Growth to 2030}
Demand is grown from 2011 at the JICA and Department of Roads 2024 main section forecast of 4.52
percent a year: by a factor of 1.94 to 2026 and 2.32 to the end of 2030, each as its own origin
destination table over the same pairs, plus a factor of 1.20 for comparison with the first
version. Every arm, baseline included, runs a six hour horizon on 15 paired seeds (notebooks 4a and
4b). The baseline completes \ProjBaseComp{} of trips; at 1.20 times demand \ProjFourComp{} complete
and time in system rises \ProjFourTis. At 2026 demand (\ProjTwentySixVeh{} vehicles) only
\ProjTwentySixComp{} of trips finish, and at 2030 demand (\ProjThirtyVeh) \ProjThirtyComp, with
time in system up \ProjThirtyTis. The network gridlocks well before 2030 demand is reached. These
levels assume no new roads, no mode shift and no peak spreading, so they mark where the present
network fails rather than forecast actual 2030 conditions.
